\documentclass{article}
\usepackage[a4paper, total={6.5in, 10in}]{geometry}
\setlength{\parindent}{0pt}
\usepackage[british]{babel}
\usepackage{amsfonts}
\usepackage{amssymb}
\usepackage{amsmath}
\usepackage{amsthm}
\usepackage{mathtools}
\usepackage{datetime2}
\usepackage{template/cenvs}

\title{\textbf{Homework 1}}
\author{Robert O'Connell}
\date{\today}

\begin{document}
\maketitle

\begin{problem}
Given points \( x,y,z,t \) in a metric space \( (X,d) \) prove that
\[
	|d(x,y) - d(z,t)| \le d(x,z) + d(y,t)
\]

\begin{solution}
	We start with the triangle inequality on \( d \)
	\begin{align*}
		d(x,y)          & \le d(x,z) + d(z,y)          \\
		d(x,y)          & \le d(x,z) + d(z,t) + d(t,y) \\
		d(x,y) - d(z,t) & \le d(x,z) + d(y,t)
	\end{align*}

	Again
	\begin{align*}
		d(z,t)          & \le d(z,x) + d(x,t)          \\
		d(z,t)          & \le d(z,x) + d(x,y) + d(y,t) \\
		d(z,t) - d(x,y) & \le d(x,z) + d(y,t)
	\end{align*}

	Then we have \( a \le b \) and \( -a \le b \). Thus \( |a| \le b \).

	i.e.
	\[
		|d(x,y) - d(z,t)| \le d(x,z) + d(y,t)
	\]
\end{solution}
\end{problem}



\begin{problem}
Suppose that \( d \) is a metric for a non-empty set \( X \), and for any \( x,y \in X \) define
\\
\( d^{(1)}(x,y) = kd(x,y) \), where \( k \) is a positive constant, \( d^{(2)}(x,y) = \min\{1,d(x,y)\} \), \( d^{(3)}(x,y) = \frac{d(x,y)}{1 + d(x,y)} \), \( d^{(4)}(x,y) = d(x,y)^2 \)
\\
Prove that \( d^{(1)}, d^{(2)}, d^{(3)} \) are metrics for \( X \) but \( d^{(4)} \) may not be a metric for \( X \)
\\

\begin{solution}
	For each \( d^{(i)}, \quad i = 1,2,3,4 \) it is clear that each function is well defined.
	Let \( x,y,z \in X \)

	Since we know that \( d \) is a metric for \( X \) it is clear that \( d^{(1)} \) preserves the symmetry and non-negativity since \( k > 0 \), also \( d^{(1)}(x,y) = 0 \iff d(x,y) = 0 \iff x = y \). Then for the triangle inequality, starting with the triangle inequality for \( d \) we have the triangle inequality of \( d^{(1)} \) once we multiply across by \( k \).
	\\

	Since \( d^{(2)} \le 0 \iff d \le  0 \), we see it preserves the non-negegativity of \( d \) and also if \( d^{(2)} = 0 \iff d = 0 \iff x=y \), also \( d^{(2)}(x,y) = \min\{1,d(x,y)\} = \min\{1, d(y,x)\} = d^{(2)}(y,x) \), so it is symmetric.

	\begin{lemma}
		For all \( a, b \ge 0\), we have \( \min\left\{ 1, a + b \right\} \le \min\left\{ 1,a \right\} + \min\left\{ 1, b \right\}  \)
	\end{lemma}
	\begin{proof}
		Case (i). \( a + b < 1 \)
		\[
			a + b \le \min\left\{ 1, a \right\} + \min\left\{ 1,b \right\} = a + b
		\]
		Since both \( a \) and \( b \) are less than one, so the min picks them out

		Case (ii) \( a + b \ge 1 \)
		\[
			1 \le \min\left\{ 1 , a \right\} + \min\left\{ 1,b \right\}
		\]
		In any selection of \( a,b \) this holds. Since if it picks out any of the ones, it is obvious, and the case condition takes care if both \( a \) and \( b \) get picked
	\end{proof}

	For the triangle inequality
	\begin{align*}
		d^{(2)}(x,y) & = \min\{1, d(x,y)\}                      \\
		             & \le \min\{1,d(x,z) + d(z,y)\}            \\
		             & \le \min\{1, d(x,z)\} + \min\{1,d(z,y)\} \\
		             & \le d^{(2)}(x,z) + d^{(2)}(z,y)
	\end{align*}

	Once again, we see that \( d^{(3)} \le 0 \iff d  \le 0 \), which is never the case, thus it is non-negative and also \( d^{(3)}(x,y) = 0 \iff d(x,y) = 0 \iff x = y  \). Also the symmetry of \( d \) implies the symmetry of \( d^{(3)} \).
	\begin{lemma}
		For \( 0 < x \le y \), we have that \( 1 - \frac{1}{x} \le 1 - \frac{1}{y} \)

		\begin{proof}
			Since \( 0 < x \le y \Rightarrow \frac{1}{x} \ge \frac{1}{y} \Rightarrow -\frac{1}{x} \le -\frac{1}{y} \Rightarrow 1 - \frac{1}{x} \le 1 - \frac{1}{y} \)
		\end{proof}
	\end{lemma}

	For the triangle inequality
	\begin{align*}
		d^{(3)}(x,y) = \frac{d(x,y)}{1+d(x,y)} & = 1- \frac{1}{1+d(x,y)}                                                     \\
		                                       & \le 1 - \frac{1}{1 + d(x,z) + d(z,y)}                                       \\
		                                       & \le \frac{d(x,z) + d(z,y)}{1 + d(x,z) + d(z,y)}                             \\
		                                       & \le \frac{d(x,z)}{1 + d(x,z) + d(z,y)} + \frac{d(z,y)}{1 + d(x,z) + d(z,y)} \\
		                                       & \le \frac{d(x,z)}{1 + d(x,z)} + \frac{d(z,y)}{1 + d(z,y)}                   \\
		                                       & \le d^{(3)}(x,z) + d^{(3)}(z,y)
	\end{align*}

	Then for \( d^{(4)} \).

	Assume it is a metric. Let \( X = \mathbb{R} \), and choose \( x = 1, z = 1.5, y = 2 \)

	Then \( d(x,y) = 1, d(x,z) = 0.5, d(z,y) = 0.5 \) with \( d \) being the usual Euclidean metric

	Looking at the triangle inequality of \( d^{(4)} \), we get
	\[
		1 \le 0.25 + 0.25
	\]
	A contradiction, thus \( d^{(4)} \) is not a metric.
\end{solution}
\end{problem}

\begin{problem}
Let \( (X,d_X) \) and \( (Y,d_Y) \) be metric spaces. For \( (x_1,y_1), (x_2,y_2) \in X \times Y \) let
\begin{align*}
	d_1((x_1,y_1), (x_2,y_2))        & = d_X(x_1,x_2) + d_Y(y_1,y_2),                \\
	d_2((x_1,y_1), (x_2,y_2))        & = \sqrt{d_X(x_1,x_2)^2 + d_Y(y_1, y_2)^2},    \\
	d_{\infty}((x_1,y_1), (x_2,y_2)) & = \max\{d_X(x_{1}, x_{2}), d_Y(y_{1},y_{2})\}
\end{align*}
\begin{itemize}
	\item Prove that for each of \( d_1,d_2,d_{\infty} \) is a metric for \( X \times Y \)
	\item Prove that for any \( p,q \in X \times Y \)
	      \[
		      d_{\infty}(p,q) \le d_2(p,q) \le d_1(p,q) \le 2d_{\infty}(p,q)
	      \]
	\item Show that the open sets in \( X \times Y \) are the same for \( d_1, d_2,d_{\infty} \)
	\item Let \( U,V \) be open subsets of \( X,Y \) respectively. Show that \( U \times V \) is \( d_i \)-open in \( X \times Y \) for \( i = 1,2, \infty \)
\end{itemize}

\begin{solution}1.

	Clearly each of the metrics \( d_1,d_{2}, d_{\infty} \) are well defined due to the properties of \( d_X \) and \( d_Y \)

	Let \( x_{1},x_{2},x_{3} \in X \) and \( y_{1}, y_{2},y_{3} \in  Y \)
	\\

	Clearly \( d_1 \) is non-negative since the sum of two non-neg is non-neg

	Also \( d_1 = 0 \iff x_{1} = x_{2} \text{ and } y_{1} = y_{2} \) implying \( (x_{1},y_{1}) = (x_{2},y_{2}) \)

	Clearly \( d_1 \) is symmetric, since both \( d_X \) and \( d_Y \) are

	Triangle Inequality
	\begin{align*}
		d_{1}((x_{1},y_{1}),(x_{2},y_{2})) & = d_X(x_{1},x_{2}) + d_Y(y_{1},y_{2})                                         \\
		                                   & \le d_X(x_{1},x_{3}) + d_X(x_{3},x_{2}) + d_Y(y_{1},y_{3}) + d_Y(y_{3},y_{2}) \\
		                                   & \le d_1((x_{1},y_{1}),(x_{3},y_{3})) + d_{1}((x_{3},y_{3}),(x_{2},y_{2}))
	\end{align*}

	Thus \( d_1 \) is a metric on \( X \times Y \)
	\\

	\( d_{2} \) is non-negative since \( \sqrt{a} \ge 0 \ \forall \ a \ge 0 \)

	Similarly \( d_{2} = 0 \iff \) the square root is 0 \( \iff d_X \text{ and } d_Y \) are zero \( \iff \) the pairs are the same

	Again due to the symmetry of \( d_X \text{ and } d_Y \) it is clear that \( d_2 \) is symmetric

	\begin{lemma}
		For \( a,b,c,d \ge 0 \), we have \( \sqrt{(a+b)^2 + (c+d)^2} \le \sqrt{a^2+c^2} + \sqrt{b^2+d^2} \)

		\begin{proof}
			Since both sides are non-negative it is equivalent to prove the following
			\[
				(a+b)^2+(c+d)^2 \le (a^2+c^2) + (b^2+d^2) + 2\sqrt{(a^2+c^2)(b^2+d^2)}
			\]
			Then expanding the LHS and cancelling out some terms we get that it is equivalent to proving
			\[
				ab + cd \le \sqrt{(a^2+c^2)(b^2+d^2)}
			\]
			Squaring and expanding
			\[
				a^2b^2+2abcd+c^2d^2 \le a^2b^2+a^2d^2+c^2b^2+c^2d^2
			\]
			\[
				0 \le a^2d^2-2abcd +c^2b^2
			\]
			\[
				0 \le (ad-bc)^2
			\]
			Which is a square and thus always non-negative.

			Hence the inequality holds
		\end{proof}
	\end{lemma}

	For the triangle inequality

	\begin{align*}
		d_{2}((x_{1},y_{1}), (x_{2},y_{2})) & = \sqrt{d_X(x_{1},x_{2})^2 + d_Y(y_{1},y_{2})^2}                                                    \\
		                                    & \le \sqrt{(d_X(x_{1},x_{3}) + d_X(x_{3},x_{2}))^2 + (d_Y(y_{1},y_{3}) + d_Y(y_{3},y_{2}))^2}        \\
		                                    & \le \sqrt{d_X(x_{1},x_{3})^2 + d_Y(y_{1},y_{3})^2} + \sqrt{d_X(x_{3},x_{2})^2 + d_Y(y_{3},y_{2})^2} \\
		                                    & \le d_{2}((x_{1},y_{1}),(x_{3},y_{3})) + d_{2}((x_{3},y_{3}),(x_{2},y_{2}))
	\end{align*}
	\\

	Again, the max of two non-negative numbers is obviously non-negative

	Also if the max of two non-neg numbers is 0, then all are 0 implying \( (x_{1},y_{1})  = (x_{2},y_{2})\)

	Again due to the symmetry of \( d_X \text{ and } d_Y \) it is clear that \( d_{\infty} \) is symmetric

	\begin{lemma}
		For all \( a,b,c,d \ge 0 \), we have \( \max\left\{ a+b, c + d \right\} \le \max\left\{ a,c \right\} + \max\left\{ b,d \right\}  \)
	\end{lemma}
	\begin{proof}
		Let \( M = \max\{a,c\} + \max\{b,d\} \)

		Since \( a \le \max\{a,c\} \) and \( b \le \max\{b,d\} \), we have \( a + b \le M \)

		Similarly \( c \le \max\{a,c\} \) and \( d \le \max\{b,d\} \) give \( c + d \le M \)

		Then \( \max\{a+b, c+d\} \le M \).
	\end{proof}

	For the triangle inequality, we get
	\begin{align*}
		d_{\infty}((x_{1},y_{1}),(x_{2},y_{2})) & = \max\left\{ d_X(x_{1},x_{2}), d_Y(y_{1},y_{2}) \right\}                                                             \\
		                                        & \le \max\left\{ d_X(x_{1},x_{3}) + d_X(x_{3},x_{2}), d_Y(y_{1},y_{3}) + d_Y(y_{3},y_{2}) \right\}                     \\
		                                        & \le \max\left\{ d_X(x_{1},x_{3}), d_Y(y_{1},y_{3}) \right\} + \max\left\{ d_X(x_{3},x_{2}), d_Y(y_{3},y_{2}) \right\} \\
		                                        & \le d_{\infty}((x_{1},y_{1}),(x_{3},y_{3})) + d_{\infty}((x_{3},y_{3}),(x_{2},y_{2}))
	\end{align*}

	Thus each of the functions are metrics on \( (X \times Y) \)

\end{solution}


\begin{solution}2.

	Let \( d_{\infty}((x_{1},y_{1}),(x_{2},y_{2})) = \max\left\{ d_X(x_1,x_{2}), d_Y(y_{1},y_{2}) \right\} = \alpha\). Then
	\begin{equation}
		d_2((x_{1},y_{1}),(x_{2},y_{2})) \ge \sqrt{\alpha ^2 + 0^2} = \alpha = d_{\infty}((x_{1},y_{1}),(x_{2},y_{2}))
	\end{equation}

	Let \( d_1((x_{1},y_{1}),(x_{2},y_{2})) = a + b, \ a,b \ge 0 \)
	Then \( d_2((x_{1},y_{1}),(x_{2},y_{2})) = \sqrt{a^2+b^2}  \)

	Then
	\[
		d_1^2 = a^2 + b^2 + 2ab, \quad d_2^2 = a^2 + b^2
	\]
	\begin{equation}
		d_2^2 \le d_{1}^2 \Rightarrow d_2((x_{1},y_{1}),(x_{2},y_{2})) \le d_1((x_{1},y_{1}),(x_{2},y_{2}))
	\end{equation}

	Again let \( d_{\infty} = \alpha \). Then \( d_1 \le \alpha + \alpha = 2d_{\infty} \)
	\begin{equation}
		\Rightarrow d_{1}((x_{1},y_{1}),(x_{2},y_{2})) \le 2 d_{\infty}((x_{1},y_{1}),(x_{2},y_{2}))
	\end{equation}

	Combining 1,2 and 3 we get, with \( p = (x_{1},y_{1}) \) and \( q = (x_{2},y_{2}) \)
	\[
		d_{\infty}(p,q) \le d_2(p,q) \le d_1(p,q) \le 2d_{\infty}(p,q)
	\]
\end{solution}

\begin{solution}3.

	Let \( U \subset X \times Y \), \( U \) is open iff
	\[
		\forall \ (x_{0},y_{0}) \in U \ \exists \ \epsilon \text{ s.t. } B_{\epsilon}(x_{0},y_{0}) \subset U
	\]
	To prove the open sets are the same under the three different metrics, we must prove that if a set is \( d_{\infty} \)-open then it is also \( d_{1} \)-open and then also \( d_{2} \)-open and finally once again to \( d_{\infty} \) completing the cycle.

	Let \( (x_{0},y_{0}) \in U \)
	\[
		B_{\epsilon_{\infty}}^{d_{\infty}}(x_{0},y_{0}) = \left\{ (x,y) \in X \times Y \mid \max\left\{ d_X(x,x_{0}), d_Y(y,y_{0}) \right\} < \epsilon_{\infty} \right\}
	\]
	\[
		B_{\epsilon_{1}}^{d_1}(x_{0},y_{0}) = \left\{ (x,y) \in X\times Y \mid d_X(x,x_{0}) + d_Y(y,y_{0}) < \epsilon_{1}  \right\}
	\]
	\[
		B_{\epsilon_{2}}^{d_{2}}(x_{0},y_{0}) = \left\{ (x,y) \in X\times Y \mid \sqrt{d_x(x,x_{0})^2 + d_Y(y,y_{0})^2} < \epsilon_{2}  \right\}
	\]
	\\

	Assume \( U \) is \( d_{\infty} \)-open. i.e. \( \exists \ \epsilon_{\infty} > 0 \text{ s.t. } B_{\epsilon_{\infty}}^{d_{\infty}}(x_{0},y_{0}) \subset U\)

	Choose \( \epsilon_{1} = \epsilon_{\infty} \)

	Then if \( (x,y) \in B_{\epsilon_{1}}^{d_1}(x_{0},y_{0}) \Rightarrow (x,y) \in B_{\epsilon_{\infty}}^{d_{\infty}}(x_{0},y_{0}) \) since \( d_{\infty} \le d_1 \) (from part above) \( \Rightarrow B_{\epsilon_{1}}^{d_{1}}(x_{0},y_{0}) \subset B_{\epsilon_{\infty}}^{d_{\infty}}(x_{0},y_{0}) \subset U  \).

	Hence \( U \) is \( d_{1} \)-open
	\\

	Assume \( U \) is \( d_{1} \)-open. i.e. \( \exists \ \epsilon_1 > 0 \text{ s.t. } B_{\epsilon_{1}}^{d_{1}}(x_{0},y_{0}) \subset U \)

	Choose \( \epsilon_2 = \frac{1}{2}\epsilon_{1} \)

	From the last part we have that \( d_1 \le 2d_{2} \)

	So if \( (x,y) \in B_{\epsilon_{2}}^{d_{2}}(x_{0},y_{0}) \) we get
	\[
		\Rightarrow d_2((x,y),(x_{0},y_{0})) < \frac{1}{2} \epsilon_{1}
	\]
	\[
		\Rightarrow d_{1}((x,y),(x_{0},y_{0})) \le 2d_{2}((x,y),(x_{0},y_{0})) < \epsilon_{1}
	\]
	So \( (x,y) \in B_{\epsilon_{1}}^{d_{1}}(x_{0},y_{0})\)

	Hence \( B_{\epsilon_{2}}^{d_{2}}(x_{0},y_{0}) \subset U \), and thus \( U \) is \( d_2 \)-open
	\\

	Assume \( U \) is \( d_2 \)-open. i.e. \( \exists \ \epsilon_{2} > 0 \text{ s.t. } B_{\epsilon_{2}}^{d_{2}}(x_{0},y_{0}) \subset U \)

	Choose \( \epsilon_{\infty} = \frac{1}{2}\epsilon_{2} \)

	Again from the last part we have that \( d_2 \le 2d_{\infty} \)

	So similar to the last time, if \( (x,y) \in  B_{\epsilon_{\infty}}^{d_{\infty}}(x_{0},y_{0}) \), we get
	\[
		d_{\infty} < \frac{1}{2} \epsilon_{2}
	\]
	\[
		\Rightarrow d_2 \le 2d_{\infty} < \epsilon_{2}
	\]

	Then \( B_{\epsilon_{\infty}}^{d_{\infty}}(x_{0},y_{0}) \subset U \) and hence \( U \) is \( d_{\infty} \)-open, completing the cycle.
\end{solution}

\begin{solution}4.
	\begin{lemma}
		Let \( A \subset B, C \subset D \)

		Then \( A\times C \subset B\times D \)

		\begin{proof}
			\begin{align*}
				A \subset B & \iff a \in A \Rightarrow a \in B \\
				C \subset D & \iff c \in C \Rightarrow c \in D
			\end{align*}
			From the properties of the Cartesian Product we get
			\[
				(x,y) \in  X\times Y \iff x \in X \text{ and } y \in Y
			\]
			Combining these we get
			\[
				(a,c) \in A \times C \Rightarrow (a,c) \in B \times D
			\]
			Therefore \( A \times C \subset B \times  D \)
		\end{proof}
	\end{lemma}

	\( U \subset X \) open and \( V \subset Y \) open

	Thus \( \forall \ x_{0} \in  U \ \exists \ \epsilon_x \text{ s.t. } B_{\epsilon_x}^{d_X}(x_{0}) \subset U \) and \( \forall \ y_{0} \in V \ \exists \ \epsilon_y \text{ s.t. } B_{\epsilon_y}^{d_Y}(y_{0}) \subset V \)

	It suffices to prove that \( U \times V \subset X \times Y \) is \( d_i \)-open for some \( i = 1,2,\infty \) as a consequence of the last part. I will choose to do it for \( d_{\infty} \)
	\\

	From Lemma 5, we have that \( U \times V \subset X \times Y \)

	\( U \times V \subset X \times Y \) is open iff
	\[
		\forall \ (x_{0},y_{0}) \in U \times V \ \exists \ \epsilon_{\infty} \text{ s.t. } B_{\epsilon_{\infty}}^{d_{\infty}}(x_{0},y_{0}) \subset U \times V
	\]
	\[
		B_{{\epsilon_{\infty}}}^{d_{\infty}}(x_{0},y_{0}) = \left\{ (x,y) \in X \times Y \mid \max\left\{ d_X(x,x_{0}),d_Y(y,y_{0}) \right\} < \epsilon_{\infty} \right\}
	\]
	Choosing \( \epsilon_{\infty} = \min\left\{ \epsilon_x,\epsilon_y \right\}  \), we have that \( d_X(x,x_{0}) < \epsilon_x \) and \( d_Y(y,y_{0}) < \epsilon_y \)

	Meaning if that \( (x,y) \in B_{\epsilon_{\infty}}^{d_{\infty}}(x_{0},y_{0}) \Rightarrow (x,y) \in B_{\epsilon_x}^{d_X}(x_{0}) \times B_{\epsilon_y}^{d_Y}(y_{0}) \)

	So then \( B_{\epsilon_{\infty}}^{d_{\infty}}(x_{0},y_{0}) \subset B_{\epsilon_x}^{d_X}(x_{0}) \times B_{\epsilon_y}^{d_Y}(y_{0}) \subset U \times V \)
	\\

	Therefore \( U \times V \) is \( d_{\infty} \)-open

	Hence it is \( d_i \)-open for \( i = 1,2,\infty \)

\end{solution}



\end{problem}

\begin{problem}
Let \( (X,d) \) be a metric space and consider \( X \times X \) as a metric space with the metric \( d_1 \) of Problem 3. Show that \( d: X \times X \to \mathbb{R}\) is continuous
\[
	d_1((x_1,x_2), (x_3,x_4)) = d(x_1,x_{3}) + d(x_{2},x_{4})
\]
\\

\begin{solution}
	Let \( (x_{0},y_{0}) \in X \times X \)

	\( d: X \times X \to \mathbb{R} \) is continuous at \( (x_{0}, y_{0}) \) iff
	\[
		\forall \ \epsilon > 0 \ \exists \ \delta > 0 \text{ s.t. } d_{1}((x_{0},y_0), (x,y)) < \delta \Rightarrow |d(x_{0},y_{0}) - d(x,y)| < \epsilon
	\]
	Let \( \epsilon > 0 \)

	Choose \( \delta = \epsilon \)

	Then \( d_1((x_{0},y_{0}), (x,y)) < \delta \) gives
	\[
		d(x_{0},x) + d(y_{0},y) < \epsilon
	\]
	Then using the property from Problem 1, we get
	\[
		|d(x_{0},y_{0}) - d(x,y)| \le d(x_{0},x) + d(y_{0},y) < \epsilon
	\]
	Which satisfies the definition of continuity above
\end{solution}
\end{problem}

\begin{problem}
Consider the set of all maps from the natural numbers to the real numbers
\[
	X = \mathbb{R}^{\mathbb{N}}
\]
Define
\[
	d(x,y) = \sum_{j=1}^{\infty} \frac{1}{2^j}\frac{|x_j-y_j|}{1+|x_j-y_j|}
\]
Show that \( (X,d) \) is a metric space

\begin{solution}
	Let \( x,y,z \in X \)

	First I should check whether or not the metric is well-defined since it is not immeaditely obvious.

	Each term in the series of partial sums of \( d(x,y) \) is non-negative since \( |x_j-y_j| \ge 0 \ \forall \ j \)

	Also, since \( \frac{|x_j-y_j|}{1 + |x_j - y_j|} \le 1 \) for all \( j \), we get
	\[
		0 \le d(x,y) \le \sum_{j=1}^{\infty} \frac{1}{2^j}
	\]

	Since, \( \sum_{j=1}^{\infty} \frac{1}{2^j} \) converges, so its limit is an upper bound on the series of partial sums, which is also monotone since it is the sum of non-negative numbers.

	Hence it is convergent and thus \( d(x,y) \) is well defined
	\\

	It is obvious that \( d(x,y) \) is non-negative since \( |x_j-y_j| \ge 0 \ \forall \ j \)

	Then if \( d(x,y)  = 0 \) we see that this happens iff \( |x_j - y_j|  = 0 \ \forall \ j\) i.e. \( x = y \)

	Also  \( |x_j - y_j| = |y_j - x_j| \) and thus \( d(x,y) = d(y,x) \)

	Then for the triangle inequality (using lemma 2 from earlier)

	\begin{align*}
		d(x,y) & = \sum_{j=1}^{\infty} \frac{1}{2^j} \frac{|x_j-y_j|}{1 + |x_j-y_j|}                                                                                \\
		       & = \sum_{j=1}^{\infty} \frac{1}{2^j}\left( 1 - \frac{1}{1 + |x_j-y_j|} \right)                                                                      \\
		       & \le \sum_{j=1}^{\infty} \frac{1}{2^j}\left( 1 - \frac{1}{1 + |x_j-z_j| + |z_j- y_j|} \right)                                                       \\
		       & \le \sum_{j=1}^{\infty} \frac{1}{2^j}\left( \frac{|x_j - z_j| + |z_j - y_j|}{1 + |x_j-z_j| + |z_j-y_j|} \right)                                    \\
		       & \le \sum_{j=1}^{\infty} \frac{1}{2^j}\left( \frac{|x_j-z_j|}{1 + |x_j-z_j| + |z_j-y_j|} + \frac{|z_j-y_j|}{1 + |x_j - z_j| + | z_j - y_j|} \right) \\
		       & \le \sum_{j=1}^{\infty} \frac{1}{2^j}\left( \frac{|x_j-z_j|}{1 + |x_j-z_j|} + \frac{|z_j-y_j|}{1 + | z_j - y_j|} \right)                           \\
		       & \le d(x,z) + d(z,y)
	\end{align*}

	Hence \( (X,d) \) is a metric space
\end{solution}
\end{problem}
\end{document}
